% ============================================================================
% chap5_parallel_orchestration.tex
%   第5章 §5.2.1 · 容器化与并行调度机制
%   从统一镜像 → 调度器 → N 并行容器 → 共享卷 → 训练器 的数据/控制链路
%
% 独立编译:
%   xelatex figs/chap5_parallel_orchestration.tex
%   → figs/chap5_parallel_orchestration.pdf
%
% 嵌入论文:
%   \begin{figure}[htbp]
%     \centering
%     \includegraphics[width=0.95\linewidth]{figs/chap5_parallel_orchestration.pdf}
%     \caption{容器化并行调度机制: 统一镜像派生 N 个隔离容器执行并行 rollout,
%              通过共享卷回收轨迹供训练器更新, 更新后的策略再分发至各容器.}
%     \label{fig:chap5_parallel_orch}
%   \end{figure}
% ============================================================================

\documentclass[tikz,border=6pt]{standalone}
\usepackage[UTF8]{ctex}
\usepackage{amsmath}
\usepackage{tikz}
\usetikzlibrary{positioning, arrows.meta, shapes.geometric, calc, fit, backgrounds}

% 配色与前两张图保持一致：橙色边框统一为深橙色
\definecolor{cOurs}{HTML}{0B4F8A}      % 深蓝：容器与策略下发
\definecolor{cOursAlt}{HTML}{C05600}   % 深橙：调度器与任务入口
\definecolor{cOursEval}{HTML}{000000}  % 黑色：数据与训练链路
\definecolor{cUpstream}{HTML}{4D4D4D}  % 深灰：基础镜像
\definecolor{cBgLayer}{HTML}{F0F0F0}
\definecolor{cBgHi}{HTML}{FFF8E7}

\begin{document}
\begin{tikzpicture}[
  font=\large\sffamily,
  every node/.style={align=center, text=black},
  % 块样式
  imgbox/.style={
    draw=cUpstream, line width=0.8pt, dashed, rounded corners=3pt,
    fill=white, text=black, inner sep=4pt,
    minimum width=5.2cm, minimum height=1.1cm,
    font=\normalsize\sffamily,
  },
  orch/.style={
    draw=cOursAlt, line width=1.1pt, rounded corners=3pt,
    fill=cBgHi, text=black, inner sep=4pt,
    minimum width=5.2cm, minimum height=1.1cm,
    font=\normalsize\bfseries\sffamily,
  },
  cbox/.style={
    draw=cOurs, line width=1pt, rounded corners=2pt,
    fill=white, inner sep=3pt,
    minimum width=2.55cm, minimum height=2.25cm,
    font=\small\bfseries\sffamily, text=black,
  },
  cinner/.style={
    draw=cOurs!60, line width=0.4pt, rounded corners=1.5pt,
    fill=cOurs!6, text=black, inner sep=1.5pt,
    minimum width=2.25cm, minimum height=0.38cm,
    font=\footnotesize\sffamily,
  },
  ctitle/.style={
    font=\small\bfseries\sffamily, text=black,
    inner sep=1pt,
  },
  storage/.style={
    draw=cOursEval, line width=1pt, rounded corners=3pt,
    fill=white, text=black, inner sep=4pt,
    minimum width=8.0cm, minimum height=1.0cm,
    font=\normalsize\sffamily,
  },
  trainer/.style={
    draw=cOursEval, line width=1pt, rounded corners=3pt,
    fill=white, text=black, inner sep=4pt,
    minimum width=5.2cm, minimum height=1.0cm,
    font=\normalsize\bfseries\sffamily,
  },
  request/.style={
    draw=black!55, line width=0.6pt, rounded corners=2pt,
    fill=white, text=black, inner sep=3pt,
    minimum width=3.2cm, minimum height=0.9cm,
    font=\small\sffamily,
  },
  % 箭头
  flow/.style={
    -{Stealth[length=2.4mm, width=2mm]},
    line width=0.9pt, draw=black!70,
  },
  flowSpawn/.style={
    -{Stealth[length=2.2mm, width=1.8mm]},
    line width=0.9pt, draw=cOursAlt,
  },
  flowData/.style={
    -{Stealth[length=2.2mm, width=1.8mm]},
    line width=0.9pt, draw=cOursEval,
  },
  flowPolicy/.style={
    -{Stealth[length=2.2mm, width=1.8mm]},
    line width=0.9pt, draw=cOurs, densely dashed,
  },
  lbl/.style={font=\footnotesize, inner sep=1pt, fill=white, text=black},
]

% ============================================================================
% 顶行: 镜像  +  用户请求
% ============================================================================
\node[imgbox] (img) at (-4.0cm, 4.8cm)
  {统一容器镜像 (Docker)\\
   \small Webots + ROS\,2 + MAPPO 训练组件};

\node[request] (req) at (4.0cm, 4.8cm)
  {用户 / 训练任务\\
   \small $K$ 组并行采样请求};

% ============================================================================
% 第 2 行: 调度器
% ============================================================================
\node[orch] (orch) at (0, 3.0cm)
  {容器运行管理 (Docker Engine)\\
   \small 创建 / 启动 / 暂停 / 销毁 · 按需分配 CPU/GPU};

\draw[flow, draw=cUpstream] (img.south) to[out=-90,in=160]
  node[lbl, pos=0.5]{构建}
  ($(orch.north west)+(1.0cm,0)$);
\draw[flow, draw=black!55] (req.south) to[out=-90,in=20]
  node[lbl, pos=0.5]{提交}
  ($(orch.north east)+(-1.0cm,0)$);

% ============================================================================
% 第 3 行: N 个并行容器 (横向排布 4 个 + 省略号)
% ============================================================================
\def\cY{0.35cm}
\node[cbox] (c1) at (-6.3cm, \cY) {};
\node[cbox] (c2) at (-3.4cm, \cY) {};
\node[cbox] (c3) at (-0.5cm, \cY) {};
\node[font=\huge\bfseries, text=black] at (1.75cm, \cY) {$\cdots$};
\node[cbox] (cN) at (4.0cm, \cY) {};

% 容器标题放在顶部 (box 内部顶对齐)
\foreach \c/\name in {c1/容器 1, c2/容器 2, c3/容器 3, cN/容器 $K$}{
  \node[ctitle, below=0.06cm of \c.north, anchor=north] (\c t) {\name};
  \node[cinner, below=0.08cm of \c t.south, anchor=north] (\c p) {MAPPO 策略模块};
  \node[cinner, below=0.08cm of \c p.south, anchor=north] (\c r) {ROS\,2 通信接口};
  \node[cinner, below=0.08cm of \c r.south, anchor=north] (\c s) {Webots 仿真};
}

% 调度器 → 容器 (spawn)
\draw[flowSpawn] ($(orch.south west)+(0.7cm,0)$)  -- (c1.north);
\draw[flowSpawn] ($(orch.south)+(-0.9cm,0)$)      -- (c2.north);
\draw[flowSpawn] ($(orch.south)+(0.9cm,0)$)       -- (c3.north);
\draw[flowSpawn] ($(orch.south east)+(-0.7cm,0)$) -- (cN.north);

\node[font=\small\bfseries, text=black, fill=white,
      draw=cOursAlt!50, rounded corners=1.5pt, inner sep=2pt]
  at ($(orch.south)+(1.8cm,-0.5cm)$) {启动 $\times K$};

% ============================================================================
% 第 4 行: 共享卷
% ============================================================================
\node[storage, minimum width=10.5cm] (vol) at (-1.15cm, -2.6cm)
  {统一数据目录 runs/\\
   \small 轨迹缓存 · 策略参数 · TensorBoard 日志};

% 各容器 → 共享卷 (数据上行)
\foreach \c in {c1,c2,c3,cN}{
  \draw[flowData] (\c.south) -- (\c.south |- vol.north);
}

\node[font=\footnotesize, text=black, fill=white, draw=cOursEval!50,
      rounded corners=1.2pt, inner sep=1.5pt]
  at ($(c3.south)+(1.45cm,-0.55cm)$) {采样轨迹与参数写入};

% ============================================================================
% 第 5 行: 训练器
% ============================================================================
\node[trainer] (tr) at (-1.15cm, -4.3cm)
  {MAPPO 训练端};

\draw[flowData] (vol.south) -- node[lbl, pos=0.5]{批量样本} (tr.north);

% 训练器 → Container K (策略下行, 蓝虚线)
\coordinate (ptop) at ($(cN.east)+(0.85cm,0)$);
\draw[flowPolicy] (tr.east) to[out=0,in=-90] (ptop) -- (cN.east);
\node[font=\small\bfseries\sffamily, text=black, align=center,
      fill=white, draw=cOurs!50, rounded corners=1.5pt, inner sep=2pt]
  at ($(ptop)+(-0.1cm,-1.9cm)$) {策略下发};

% ============================================================================
% 图例
% ============================================================================
\node[draw=black!25, rounded corners=2pt, inner sep=5pt, font=\small,
      text=black, anchor=north]
  at ($(tr.south)+(0,-0.35cm)$)
  {\begin{tabular}{@{}l@{\hspace{4pt}}l@{\hspace{14pt}}l@{\hspace{4pt}}l@{\hspace{14pt}}l@{\hspace{4pt}}l@{\hspace{14pt}}l@{\hspace{4pt}}l@{}}
     \tikz[baseline=-0.6ex]\draw[flowSpawn] (0,0) -- (0.6,0); & 容器启动 &
     \tikz[baseline=-0.6ex]\draw[flowData] (0,0) -- (0.6,0); & 数据上行 (采样轨迹) &
     \tikz[baseline=-0.6ex]\draw[flowPolicy] (0,0) -- (0.6,0); & 策略下行 (参数更新) &
     \tikz[baseline=-0.6ex]\draw[flow, draw=black!55] (0,0) -- (0.6,0); & 用户请求 / 构建 \\
   \end{tabular}};

\end{tikzpicture}
\end{document}